% ch03_controle.tex — Structures de Contr\^ole
\chapter{Structures de Contr\^ole --- Conditions et Boucles}
\label{ch:controle}

\begin{intuition}
Imaginez une recette de cuisine~: <<~Si la p\^ate est trop liquide, ajoutez de la farine~;
sinon, enfournez pendant 30 minutes.~>> Un programme fonctionne exactement de la m\^eme
mani\`ere~: il prend des \emph{d\'ecisions} et \emph{r\'ep\`ete} des \'etapes selon des conditions.
Les structures de contr\^ole sont le c\oe ur de tout algorithme scientifique.
\end{intuition}

% ============================================================
\section{Les conditions~: \py{if}, \py{elif}, \py{else}}
\index{if}\index{elif}\index{else}\index{condition}

En science, on classifie souvent selon des seuils. Consid\'erons les
\emph{transitions de phase}\index{transition de phase} de l'eau~:

\begin{codeblock}[title=\textbf{Python}]
\begin{minted}{python}
temperature = -5  # en degrés Celsius

if temperature < 0:
    etat = "solide"
elif temperature < 100:
    etat = "liquide"
else:
    etat = "gazeux"

print(f"À {temperature}°C, l'eau est {etat}.")
\end{minted}
\end{codeblock}

\begin{codeoutput}
\begin{verbatim}
À -5°C, l'eau est solide.
\end{verbatim}
\end{codeoutput}

\begin{remarque}
Python \'evalue les conditions \emph{dans l'ordre}. D\`es qu'une condition est vraie,
le bloc correspondant est ex\'ecut\'e et les suivants sont ignor\'es.
\end{remarque}

% --- Flowchart TikZ ---
\begin{figure}[ht]
\centering
\begin{tikzpicture}[
    decision/.style={diamond, draw, fill=blue!10, text width=3cm,
        align=center, inner sep=2pt},
    block/.style={rectangle, draw, fill=green!10, rounded corners,
        minimum height=1cm, text width=3cm, align=center},
    arr/.style={->, thick, >=stealth}
]
\node[block] (start) {D\'ebut};
\node[decision, below=1cm of start] (d1) {$T < 0$~?};
\node[block, right=2.5cm of d1] (sol) {\'Etat = solide};
\node[decision, below=1.5cm of d1] (d2) {$T < 100$~?};
\node[block, right=2.5cm of d2] (liq) {\'Etat = liquide};
\node[block, below=1.5cm of d2] (gaz) {\'Etat = gazeux};

\draw[arr] (start) -- (d1);
\draw[arr] (d1) -- node[above]{oui} (sol);
\draw[arr] (d1) -- node[left]{non} (d2);
\draw[arr] (d2) -- node[above]{oui} (liq);
\draw[arr] (d2) -- node[left]{non} (gaz);
\end{tikzpicture}
\caption{Organigramme d'une structure \py{if/elif/else} pour les transitions de phase.}
\label{fig:flowchart-if}
\end{figure}

\subsection{Op\'erateurs de comparaison et op\'erateurs logiques}
\index{op\'erateur de comparaison}\index{op\'erateur logique}

\begin{definition}[Op\'erateurs de comparaison]
Les op\'erateurs \py{==}, \py{!=}, \py{<}, \py{>}, \py{<=}, \py{>=} renvoient
un bool\'een (\py{True} ou \py{False}). On peut combiner des conditions avec
\py{and}, \py{or}, \py{not}.
\end{definition}

\begin{codeblock}[title=\textbf{Python}]
\begin{minted}{python}
pH = 3.2

if pH < 7 and pH >= 0:
    print("Solution acide")
elif pH == 7:
    print("Solution neutre")
elif pH > 7 and pH <= 14:
    print("Solution basique")
else:
    print("Valeur de pH invalide")
\end{minted}
\end{codeblock}

\begin{codeoutput}
\begin{verbatim}
Solution acide
\end{verbatim}
\end{codeoutput}

% ============================================================
\section{La boucle \py{while}}
\index{while}\index{boucle}

La boucle \py{while} r\'ep\`ete un bloc \emph{tant qu'une condition est vraie}.
En sciences, on l'utilise souvent pour des algorithmes de \emph{convergence}\index{convergence}.

\begin{exemple}[Convergence d'une suite]
La suite $u_{n+1} = \frac{1}{2}\left(u_n + \frac{a}{u_n}\right)$ converge vers $\sqrt{a}$
(m\'ethode de H\'eron)\index{m\'ethode de H\'eron}.
\end{exemple}

\begin{codeblock}[title=\textbf{Python}]
\begin{minted}{python}
a = 2.0
u = 1.0       # valeur initiale
tolerance = 1e-10
iterations = 0

while abs(u**2 - a) > tolerance:
    u = 0.5 * (u + a / u)
    iterations += 1

print(f"sqrt({a}) ≈ {u}")
print(f"Convergence en {iterations} itérations")
\end{minted}
\end{codeblock}

\begin{codeoutput}
\begin{verbatim}
sqrt(2.0) ≈ 1.4142135623730951
Convergence en 4 itérations
\end{verbatim}
\end{codeoutput}

\begin{attention}
Une boucle \py{while} dont la condition ne devient \emph{jamais} fausse tourne
\`a l'infini~! Ajoutez toujours un garde-fou~:
\begin{minted}{python}
max_iter = 10000
while condition and iterations < max_iter:
    ...
\end{minted}
\end{attention}

% ============================================================
\section{La boucle \py{for} et \py{range()}}
\index{for}\index{range}

La boucle \py{for} parcourt un \emph{it\'erable} (liste, cha\^ine, \py{range}, etc.).

\begin{codeblock}[title=\textbf{Python}]
\begin{minted}{python}
# Calculer la somme des carrés de 1 à 10
somme = 0
for i in range(1, 11):
    somme += i**2

print(f"Somme des carrés de 1 à 10 = {somme}")
# Vérification : n(n+1)(2n+1)/6
print(f"Formule : {10 * 11 * 21 // 6}")
\end{minted}
\end{codeblock}

\begin{codeoutput}
\begin{verbatim}
Somme des carrés de 1 à 10 = 385
Formule : 385
\end{verbatim}
\end{codeoutput}

\begin{definition}[\py{range(start, stop, step)}]
\py{range(a, b, s)} g\'en\`ere les entiers de \py{a} \`a \py{b-1} avec un pas de \py{s}.
Attention~: la borne sup\'erieure \py{b} est \emph{exclue}.
\end{definition}

\begin{codeblock}[title=\textbf{Python}]
\begin{minted}{python}
# Afficher les multiples de 3 entre 0 et 15
for n in range(0, 16, 3):
    print(n, end=" ")
\end{minted}
\end{codeblock}

\begin{codeoutput}
\begin{verbatim}
0 3 6 9 12 15
\end{verbatim}
\end{codeoutput}

% ============================================================
\section{Boucles imbriqu\'ees}
\index{boucle imbriqu\'ee}

\begin{codeblock}[title=\textbf{Python}]
\begin{minted}{python}
# Table de multiplication (extrait)
for i in range(1, 4):
    for j in range(1, 4):
        print(f"{i} × {j} = {i*j:2d}", end="   ")
    print()  # retour à la ligne
\end{minted}
\end{codeblock}

\begin{codeoutput}
\begin{verbatim}
1 × 1 =  1   1 × 2 =  2   1 × 3 =  3
2 × 1 =  2   2 × 2 =  4   2 × 3 =  6
3 × 1 =  3   3 × 2 =  6   3 × 3 =  9
\end{verbatim}
\end{codeoutput}

% ============================================================
\section{\py{break} et \py{continue}}
\index{break}\index{continue}

\begin{codeblock}[title=\textbf{Python}]
\begin{minted}{python}
# Trouver le premier nombre premier > 50
n = 51
while True:
    est_premier = True
    for d in range(2, int(n**0.5) + 1):
        if n % d == 0:
            est_premier = False
            break       # inutile de tester d'autres diviseurs
    if est_premier:
        print(f"Premier nombre premier > 50 : {n}")
        break           # sortir de la boucle while
    n += 1
\end{minted}
\end{codeblock}

\begin{codeoutput}
\begin{verbatim}
Premier nombre premier > 50 : 53
\end{verbatim}
\end{codeoutput}

\begin{codeblock}[title=\textbf{Python}]
\begin{minted}{python}
# Sauter les valeurs négatives dans des mesures
mesures = [2.3, -1.0, 4.5, -0.3, 3.8, 7.1]
somme = 0
compteur = 0
for m in mesures:
    if m < 0:
        continue   # passer à l'itération suivante
    somme += m
    compteur += 1

print(f"Moyenne (valeurs positives) = {somme/compteur:.2f}")
\end{minted}
\end{codeblock}

\begin{codeoutput}
\begin{verbatim}
Moyenne (valeurs positives) = 4.43
\end{verbatim}
\end{codeoutput}

% ============================================================
\section{Erreurs courantes}

\begin{errmark}
\begin{enumerate}
    \item \textbf{Erreur d'indentation}~: Python utilise l'indentation pour d\'elimiter
    les blocs. M\'elanger espaces et tabulations provoque une \py{IndentationError}.
    R\`egle~: utilisez toujours \textbf{4 espaces}.

    \item \textbf{Boucle infinie}~: oublier d'incr\'ementer la variable de contr\^ole
    dans un \py{while} m\`ene \`a une boucle qui ne s'arr\^ete jamais.
    \begin{minted}{python}
# ERREUR : i ne change jamais !
i = 0
while i < 10:
    print(i)    # affiche 0 indéfiniment
    \end{minted}

    \item \textbf{Erreur de d\'ecalage (\emph{off-by-one})}~: \py{range(1, 10)} produit
    les entiers de 1 \`a \textbf{9}, pas 10~! Pour inclure 10, \'ecrivez \py{range(1, 11)}.

    \item \textbf{Utiliser \py{=} au lieu de \py{==}}~: \py{=} est l'affectation,
    \py{==} est la comparaison. \'Ecrire \py{if x = 5} provoque une \py{SyntaxError}.
\end{enumerate}
\end{errmark}

% ============================================================
\section{Mini-projet~: D\'esint\'egration radioactive}

\begin{projet}
La loi de d\'esint\'egration radioactive\index{d\'esint\'egration radioactive} donne le nombre
d'atomes restants~:
\[
    N(t) = N_0 \times 0{,}5^{\,t / t_{1/2}}
\]
o\`u $N_0$ est le nombre initial d'atomes et $t_{1/2}$ la demi-vie.

\textbf{Objectif}~: simuler la d\'esint\'egration du carbone-14 ($t_{1/2} = 5730$ ans)
et afficher le nombre d'atomes restants chaque 1000 ans.
\end{projet}

\begin{codeblock}[title=\textbf{Python}]
\begin{minted}{python}
N0 = 1_000_000       # nombre initial d'atomes
demi_vie = 5730      # demi-vie du C-14 en années
seuil = N0 * 0.01    # arrêter quand il reste < 1%

t = 0
N = N0
print(f"{'Temps (ans)':>12} {'N restants':>12} {'% restant':>10}")
print("-" * 36)

while N > seuil:
    pourcentage = N / N0 * 100
    print(f"{t:>12} {N:>12.0f} {pourcentage:>9.1f}%")
    t += 1000
    N = N0 * 0.5 ** (t / demi_vie)

print(f"\nAprès {t} ans, il reste moins de 1% des atomes.")
\end{minted}
\end{codeblock}

\begin{codeoutput}
\begin{verbatim}
 Temps (ans)   N restants  % restant
------------------------------------
           0      1000000     100.0%
        1000       886076      88.6%
        2000       785128      78.5%
        3000       695685      69.6%
        4000       616441      61.6%
        5000       546274      54.6%
        ...
       37000        12055       1.2%

Après 38000 ans, il reste moins de 1% des atomes.
\end{verbatim}
\end{codeoutput}

% ============================================================
\section{Exercices}

\begin{exercice}[$\star$ --- Classification de l'IMC]
\'Ecrivez un programme qui demande le poids (kg) et la taille (m) d'une personne,
calcule l'IMC ($\text{IMC} = \text{poids} / \text{taille}^2$) et affiche la
cat\'egorie~: sous-poids ($<18{,}5$), normal ($18{,}5$--$25$), surpoids ($25$--$30$),
ob\'esit\'e ($\geq 30$).
\end{exercice}

\begin{exercice}[$\star$ --- Compte \`a rebours]
Utilisez une boucle \py{while} pour afficher un compte \`a rebours de 10 \`a 0,
puis affichez <<~D\'ecollage~!~>>.
\end{exercice}

\begin{exercice}[$\star\star$ --- Suite de Fibonacci]
Calculez les 20 premiers termes de la suite de Fibonacci\index{Fibonacci}
($F_0 = 0$, $F_1 = 1$, $F_{n} = F_{n-1} + F_{n-2}$) avec une boucle \py{for}.
V\'erifiez que le rapport $F_n / F_{n-1}$ converge vers le nombre d'or
$\varphi \approx 1{,}618$.
\end{exercice}

\begin{exercice}[$\star\star$ --- Crible d'\'Eratosth\`ene]
Utilisez des boucles imbriqu\'ees pour trouver tous les nombres premiers
inf\'erieurs \`a 100.
\end{exercice}

\begin{exercice}[$\star\star\star$ --- Simulation de marche al\'eatoire]
Simulez une marche al\'eatoire 1D\index{marche al\'eatoire}~: \`a chaque pas, la particule se
d\'eplace de $+1$ ou $-1$ avec probabilit\'e \'egale. Apr\`es $N = 1000$ pas,
affichez la position finale. R\'ep\'etez 100 fois et calculez la distance
moyenne $\langle |x| \rangle$. Comparez avec la pr\'ediction th\'eorique
$\sqrt{N} \approx 31{,}6$.
\emph{Indice~:} utilisez \py{import random} et \py{random.choice([-1, 1])}.
\end{exercice}

% ============================================================
\begin{formules}{R\'esum\'e du chapitre~: Structures de contr\^ole}
\begin{itemize}
    \item \textbf{Conditions}~: \py{if condition:} / \py{elif condition:} / \py{else:}
    \item \textbf{Boucle \py{while}}~: r\'ep\`ete tant que la condition est vraie
    \item \textbf{Boucle \py{for}}~: parcourt un it\'erable (\py{range()}, liste, etc.)
    \item \textbf{\py{range(a, b, s)}}~: entiers de \py{a} \`a \py{b-1}, pas de \py{s}
    \item \textbf{\py{break}}~: sort imm\'ediatement de la boucle
    \item \textbf{\py{continue}}~: passe \`a l'it\'eration suivante
    \item \textbf{Indentation}~: 4 espaces, toujours coh\'erente
    \item \textbf{Garde-fou}~: pr\'evoir un nombre maximal d'it\'erations pour \py{while}
\end{itemize}
\end{formules}
